% !TEX root = ../main.tex

\section{Tokens}

\begin{frame}{An LLM is a next-token predictor}
  \small
  A generative model learns a probability distribution over what could come
  next. For a language model, the possible outputs are tokens from vocabulary
  $V$.

  \vspace{0.5em}
  \begin{center}
    \begin{tikzpicture}[
      stage/.style={draw, rounded corners=3pt, minimum height=0.9cm,
        align=center},
      >=Latex
    ]
      \node[stage, fill=tokenblue, draw=tokenblueborder,
        minimum width=3.7cm] (context) {\texttt{The cat sat on the}};
      \node[stage, fill=tokenorange, draw=tokenorangeborder,
        right=0.8cm of context, minimum width=3.3cm] (model)
        {language model};
      \node[stage, fill=tokengreen, draw=tokengreenborder,
        right=0.8cm of model, minimum width=3.0cm] (next)
        {next-token\\probabilities};
      \draw[->, thick] (context) -- (model);
      \draw[->, thick] (model) -- (next);
    \end{tikzpicture}
  \end{center}

  \vspace{0.3em}
  \begin{columns}[T]
    \begin{column}{0.46\textwidth}
      \centering
      \begin{tabular}{lc}
        \textbf{Candidate} & \textbf{Probability} \\
        \hline
        \tokenbox[tokengreen]{mat}   & $0.46$ \\
        \tokenbox[tokengreen]{floor} & $0.21$ \\
        \tokenbox[tokengreen]{sofa}  & $0.12$ \\
        other tokens                 & $0.21$
      \end{tabular}
    \end{column}
    \begin{column}{0.50\textwidth}
      \begin{enumerate}
        \item Choose or sample one next token.
        \item Append it to the existing sequence.
        \item Feed the longer sequence back into the model.
        \item Repeat until generation stops.
      \end{enumerate}
    \end{column}
  \end{columns}

  \vfill
  \begin{alertblock}{Why begin with tokens?}
    The model generates one token at a time, so we first need to understand
    what tokens are and how text becomes a token sequence.
  \end{alertblock}
\end{frame}

\begin{frame}{A token is a discrete piece of data}
  \begin{center}
    \begin{tikzpicture}[
      item/.style={draw, rounded corners=3pt, minimum width=3.2cm,
        minimum height=2.4cm, align=center, line width=0.8pt},
      label/.style={font=\bfseries}
    ]
      \node[item, fill=tokenblue, draw=tokenblueborder] (text) {
        \textbf{Text}\\[0.4em]
        \texttt{ing}\qquad\texttt{,}\\[0.35em]
        {\scriptsize subword or punctuation}
      };
      \node[item, fill=tokenorange, draw=tokenorangeborder,
        right=0.5cm of text] (image) {
        \textbf{Image}\\[0.35em]
        \begin{tikzpicture}[scale=0.25]
          \foreach \x in {0,...,3}
            \foreach \y in {0,...,3}
              \draw[fill=white] (\x,\y) rectangle +(1,1);
          \fill[tokenorangeborder!65] (1,1) rectangle +(2,2);
          \draw[line width=1pt] (1,1) rectangle +(2,2);
        \end{tikzpicture}\\[-0.2em]
        {\scriptsize patch or visual code}
      };
      \node[item, fill=tokengreen, draw=tokengreenborder,
        right=0.5cm of image] (audio) {
        \textbf{Audio}\\[0.35em]
        \begin{tikzpicture}[x=0.18cm,y=0.18cm]
          \draw[tokengreenborder, line width=1pt]
            plot[smooth] coordinates {(0,2) (1,3) (2,0) (3,4)
              (4,1) (5,3) (6,2) (7,2)};
        \end{tikzpicture}\\[0.25em]
        {\scriptsize short segment or codec code}
      };
    \end{tikzpicture}
  \end{center}

  \begin{block}{Common idea}
    A model processes a sequence of numbered units. The tokenizer decides
    where one unit ends and the next begins.
  \end{block}
\end{frame}

\begin{frame}{Text can be split in different ways}
  \small
  Consider the same input: \texttt{unbelievable!}

  \vspace{0.8em}
  \begin{tabular}{@{}p{2.2cm}p{8.8cm}@{}}
    \textbf{Characters} &
      \tokenbox{u}\,\tokenbox{n}\,\tokenbox{b}\,\tokenbox{e}\,
      \tokenbox{l}\,\tokenbox{i}\,\tokenbox{e}\,\tokenbox{v}\,
      \tokenbox{a}\,\tokenbox{b}\,\tokenbox{l}\,\tokenbox{e}\,
      \tokenbox{!} \\[0.8em]
    \textbf{Words} &
      \tokenbox[tokenorange]{unbelievable}\,\tokenbox[tokenorange]{!}
      \\[0.8em]
    \textbf{Subwords} &
      \tokenbox[tokengreen]{un}\,\tokenbox[tokengreen]{believ}\,
      \tokenbox[tokengreen]{able}\,\tokenbox[tokengreen]{!}
  \end{tabular}

  \vfill
  \begin{alertblock}{Why subwords?}
    Frequent strings can become one token, while rare words can still be
    assembled from smaller pieces. A token is therefore not necessarily a
    character or a complete word.
  \end{alertblock}

  {\scriptsize The subword split above is illustrative; the result depends on
  the tokenizer and its training data.}
\end{frame}

\begin{frame}{How repeated strings become tokens: toy BPE}
  \footnotesize
  \textbf{Byte Pair Encoding (BPE)} starts with characters; each word in this
  toy corpus occurs once.

  \vspace{0.15em}
  {\centering
    \tokenbox{play}\quad
    \tokenbox{played}\quad
    \tokenbox{player}\quad
    \tokenbox{playing}\par
  }

  At every round, count \emph{adjacent pairs across the whole corpus}, then
  replace one most-frequent pair with a new symbol.

  \vspace{0.2em}
  \begin{columns}[T]
    \begin{column}{0.59\textwidth}
      \centering
      {\scriptsize
      \begin{tabular}{clcl}
        \textbf{Round} & \textbf{Selected pair} & \textbf{Count} &
          \textbf{Common prefix} \\
        \hline
        0 & --                 & -- & \texttt{p l a y} \\
        1 & \texttt{p + l}    & 4  & \texttt{pl a y} \\
        2 & \texttt{pl + a}   & 4  & \texttt{pla y} \\
        3 & \texttt{pla + y}  & 4  & \texttt{play} \\
      \end{tabular}
      }

      \vspace{0.4em}
      {\scriptsize Initially several pairs tie at count 4; a fixed tie-breaking
      rule chooses which equally frequent pair is merged first.}
    \end{column}
    \begin{column}{0.37\textwidth}
      \begin{block}{What does count 4 mean?}
        The selected pair appears once in each of the four words, for four
        occurrences across the corpus.
      \end{block}

      \centering
      After round 3:\par\smallskip
      {\scriptsize
      \begin{tabular}{rcl}
        \texttt{played} & $\rightarrow$ &
          \texttt{[play][e][d]} \\
        \texttt{player} & $\rightarrow$ &
          \texttt{[play][e][r]}
      \end{tabular}
      }
    \end{column}
  \end{columns}

  \vfill
  \begin{alertblock}{Stopping criterion}
    This example stops after round 3 only to highlight \texttt{play}. Real BPE
    continues to a merge budget or target vocabulary size. Here, the next merge
    is \texttt{play + e}, which occurs twice.
  \end{alertblock}
\end{frame}

\begin{frame}{A vocabulary maps tokens to integer IDs}
  \begin{columns}[T]
    \begin{column}{0.43\textwidth}
      \centering
      \textbf{Tiny illustrative vocabulary}

      \vspace{0.5em}
      \begin{tabular}{rcl}
        \textbf{ID} & & \textbf{Token} \\
        \hline
        7   & $\leftrightarrow$ & \tokenbox[tokenorange]{!} \\
        41  & $\leftrightarrow$ & \tokenbox[tokenorange]{<space>} \\
        88  & $\leftrightarrow$ & \tokenbox[tokenorange]{er} \\
        314 & $\leftrightarrow$ & \tokenbox[tokenorange]{play} \\
      \end{tabular}
    \end{column}
    \begin{column}{0.53\textwidth}
      Input begins with a space: \texttt{player!}
      \[
        \tokenbox{<space>}\,
        \tokenbox{play}\,
        \tokenbox{er}\,
        \tokenbox{!}
      \]
      \[
        \longrightarrow [41,\,314,\,88,\,7]
      \]

      \begin{block}{IDs are labels, not quantities}
        ID 314 is not ``larger in meaning'' than ID 88. The numbers only index
        entries in a fixed list.
      \end{block}
    \end{column}
  \end{columns}

  \vfill
  \small Whitespace and punctuation also carry information. Depending on the
  tokenizer, a space may be separate or attached to a neighboring token.

  \medskip
  \begin{center}
    \textbf{The complete collection of all available tokens is the vocabulary,
    denoted $V$.}
  \end{center}
\end{frame}
